\section{Conclusiones}

El objetivo general de este trabajo era evaluar en qué medida un pipeline RAG aplicado a la búsqueda de empleo remoto mejora la calidad de la recuperación frente a la búsqueda por palabras clave, e identificar el aporte y el costo de cada componente, con JobAI como caso de estudio. En JobAI, la configuración completa eleva el nDCG@10 de 0,570 con BM25 a 0,834 ($+0{,}263$ [$+0{,}207$; $+0{,}321$]), y la búsqueda por subcadena que ofrecía la plataforma no devolvía resultados en 78 de 85 consultas. La mejora proviene de la recuperación densa, los filtros estructurados y el reordenamiento; la búsqueda híbrida no aporta. El costo se concentra en el reordenamiento, cerca de 0,8~s por consulta. Las conclusiones por objetivo específico son las siguientes.

\paragraph{Objetivo 1: limitaciones de la búsqueda léxica y componentes de RAG.} La literatura identifica en la búsqueda de empleo una limitación estructural de la búsqueda léxica, el desajuste de vocabulario entre cómo se describe una necesidad y cómo se redacta una oferta, y propone como respuesta la recuperación densa, los filtros estructurados, la búsqueda híbrida, el reordenamiento con \textit{cross-encoders} y los modelos multilingües. Cada componente tiene respaldo propio, pero medido sobre todo en colecciones de propósito general, y BM25 sigue siendo una línea base robusta que un sistema denso no siempre supera \citep{Thakur2021}. Esa revisión fijó la comparación principal del trabajo, BM25 y no la búsqueda de la plataforma, y las categorías de consulta que permiten ver dónde actúa cada componente.

\paragraph{Objetivo 2: protocolo de evaluación reproducible.} El protocolo combina un corpus real congelado de 1.068 ofertas, 100 consultas de prueba en cinco categorías, en inglés y en español, un conjunto de desarrollo separado para las decisiones de configuración, nueve sistemas que forman una escalera de componentes, métricas por necesidad de información con pruebas de permutación y corrección de Holm, y un \textit{benchmark} de almacenes vectoriales con datos sintéticos calibrados. Su elemento más delicado es el juez de relevancia. El juez final se eligió con una regla fijada antes de anotar, y su acuerdo con las etiquetas humanas en 298 pares de validación fue sustancial ($\kappa$ binario de 0,786 [0,706; 0,858]), por encima del umbral previsto. El juez que se usó durante el desarrollo, un modelo local pequeño, obtuvo 0,508 en la misma muestra: un juez automático solo es utilizable después de validarlo. El corpus, las consultas, los juicios, las corridas y el código del arnés están en el repositorio público del trabajo (Anexo~\ref{anx:repositorio}).

\paragraph{Objetivo 3: dos iteraciones de diseño y evaluación.} La Iteración~1 corrigió el almacenamiento vectorial y defectos del artefacto; la Iteración~2 agregó los filtros, la búsqueda híbrida, el reordenamiento, el modelo multilingüe y la compuerta de sin resultados. La aplicación del protocolo reveló 20 fallas: 12 del artefacto, de las cuales 10 eran silenciosas, y 8 del propio arnés de evaluación. Ninguna de las fallas silenciosas producía un error visible, porque el modelo de lenguaje convertía en una respuesta fluida lo que recibía, aunque fuera incompleto, estuviera mal filtrado o mal enrutado. Se detectaron porque el protocolo mide salidas intermedias de cada etapa, no solo la respuesta final. Dos quedan abiertas a propósito, para no ajustar el artefacto al conjunto de prueba.

\paragraph{Objetivo 4: aporte y costo de cada componente.}
\begin{itemize}
    \item \textbf{Recuperación densa (PE1).} Supera a BM25 en $+0{,}122$ de nDCG@10, en las cuatro categorías y con la mayor ventaja en las consultas de desajuste de vocabulario, que siguen siendo las más difíciles. Con el modelo multilingüe, la configuración completa en español (0,778) supera a BM25 incluso sobre la consulta en inglés (0,626).
    \item \textbf{Filtros (PE1).} Aportan $+0{,}046$ al nDCG@10 global y $+0{,}181$ en las consultas con restricciones, y elevan el cumplimiento de restricciones en las cinco primeras posiciones de 0,41 a 0,95.
    \item \textbf{Búsqueda híbrida (PE1).} No aporta ($-0{,}003$, sin diferencia significativa): mejora las consultas ruidosas y empeora las de restricciones y las de español. Con RRF sin ponderar, una lista léxica mucho más débil que la densa no se aprovecha.
    \item \textbf{Reordenamiento (PE1 y PE3).} Es el componente que más aporta ($+0{,}099$) y el que más cuesta: 811 de los 835~ms de la construcción del contexto. Con la clasificación de la consulta, el usuario espera 1,96~s en la mediana hasta el primer \textit{token}.
    \item \textbf{Fidelidad (PE2).} Con el juez final, en promedio el 90\,\% de las afirmaciones de cada respuesta está sostenido por las ofertas recuperadas (fidelidad media de 0,901), y ninguna respuesta queda por debajo de la mitad. La verificación humana de una muestra de 20 respuestas dio una fidelidad mayor (0,980), de modo que la cifra del juez es conservadora. Las pocas afirmaciones sin respaldo son sobre todo introducciones que presentan todas las ofertas como si cumplieran el criterio de la búsqueda.
    \item \textbf{Motor vectorial (PE4).} A la escala de JobAI, pgvector y Qdrant producen la misma calidad, y su diferencia de latencia desaparece frente al reordenamiento. La elección importa con filtros selectivos a gran escala: con 500.000 vectores y un filtro del 1\,\%, Qdrant mantiene un \textit{recall} de 0,982, y pgvector se estanca en 0,876 con barrido iterativo.
\end{itemize}

En síntesis, para un agregador de ofertas pequeño, con componentes de uso general y sin datos de entrenamiento propios, un pipeline RAG mejora sustancialmente la búsqueda frente a BM25, siempre que se trate como un conjunto de componentes que se miden por separado: no todos aportan, el que más aporta es también el más costoso, y varias de sus fallas solo aparecen cuando se mide cada etapa.

\section{Recomendaciones y trabajo futuro}

\begin{itemize}
    \item \textbf{Evaluación con usuarios.} La relevancia juzgada no equivale a la utilidad para quien busca empleo. Un estudio con usuarios diría si la mejora medida se traduce en búsquedas más exitosas y si la espera de unos dos segundos es aceptable.
    \item \textbf{Fusión híbrida ponderada.} Una combinación convexa de puntajes o un RRF con peso o $k$ calibrados en el conjunto de desarrollo \citep{Bruch2023}, y BM25 sobre una traducción automática de la consulta, son las alternativas naturales al resultado nulo de la búsqueda híbrida.
    \item \textbf{Reordenamiento y clasificación más baratos.} Reordenar menos candidatas, o solo cuando la primera etapa es incierta, y resolver con reglas o con un modelo más pequeño las consultas simples que hoy pasan por el clasificador basado en LLM.
    \item \textbf{Fallas abiertas.} El patrón de \textit{hiring} en el enrutamiento (A8), las plantillas de agregación faltantes (C2) y la extracción del tipo de contrato \enquote{prácticas}, que deben corregirse y evaluarse con consultas nuevas, no con las de prueba.
    \item \textbf{Normalización de salarios.} Convertir monedas y periodos de pago en lugar de descartar los valores implausibles (A5).
    \item \textbf{Diálogo de varios turnos y personalización.} Llevar el historial de la conversación a la recuperación y evaluar aparte la personalización por currículo, que en este trabajo solo se corrigió (A1).
    \item \textbf{Otros corpus.} Repetir la evaluación con ofertas en español y con un corpus más grande y de más fuentes, para separar lo que depende de JobAI de lo que es propio de la búsqueda de empleo.
\end{itemize}
